\section{Introduction}

The fundamental characteristic of finite-resolution observation is irreversibility: the observation output is a compression and coarse-graining of the ontological state; even if the ontological evolution is reversible in an information-theoretic sense, the observation layer may still exhibit multi-valuedness, irreversibility, periodization artifacts, and boundary effects. This paper aims to establish, at a strictly auditable mathematical level, a complete logical chain so that the following objects and propositions can be defined, proved, and experimentally verified within a single formal system:
\begin{enumerate}
  \item \textbf{Observation kernel:} formalizing finite observation as an irreversible canonical operator.
  \item \textbf{Stable visible states:} characterizing ``visible stable structures'' as fixed-point sets of the observation kernel.
  \item \textbf{Invisible degrees of freedom (time fibers):} rigorously objectifying the non-injectivity of the observation kernel as fibers.
  \item \textbf{Reality closure and self-reference:} defining ``self-consistent reality'' as a fixed point of the composition operator of closure and evolution.
  \item \textbf{Reversible ontological dynamics:} constructing reversible microscopic time evolution and demonstrating its macroscopic multi-valuedness under projection.
\end{enumerate}

This paper adopts a ``definition—theorem—proof—reproducible experiment—discussion'' organizational structure. The interpretive layer (e.g., high-dimensional periodic ontology and low-dimensional projection) serves only as motivation and interfacing, not feeding back into the theorem layer derivation.
